\documentclass[11pt]{beamer}
\usetheme{Warsaw}
\usepackage[utf8]{inputenc}
\usepackage[english]{babel}
\usepackage{amsmath}
\usepackage{amsfonts}
\usepackage{amssymb}
\usepackage{graphicx}
\author{Andre Lelis and Renan Andrade}
\title{Peace in an unequal world?}
%\setbeamercovered{transparent} 
%\setbeamertemplate{navigation symbols}{} 
%\logo{} 
%\institute{EPGE - FGV} 
%\subject{} 
\begin{document}

\begin{frame}
\titlepage
\end{frame}

%\begin{frame}
%\tableofcontents
%\end{frame}

\begin{frame}{Paper}

"Peace in an unequal world? Experimental evidence on the relationship between inequality and conflict in a guns-vs-butter setting"

Games and Economic Behavior,
Volume 147,2024,Pages 74-87\\

Authors: Alexandra Baier, Sophia Seelos, Thomas Rittmannsberger
\end{frame}

\begin{frame}{Motivation}

\begin{itemize}
\item Many governments spend significant portion of the income on military and security forces.
\item These investments are often sunk costs. These resources could be used in other essential social and economics needs.
\item Is unarmed peace feasible?
\end{itemize}

\end{frame}

\begin{frame}{Model}

\begin{itemize}
\item Garfinkel and Syropoulos (2021)

\item Two players ($i= A,B$). Each player divides his resources $R_{i}$ between $X_{i}$(production) and $G_{i}$(arming) such that  $R_{i} = X_{i} + G_{i}$. 

\item If neither player attack, peace is sustained and both consumer their own production

\item If at least one attack, then the sum of production constitutes the contestable output $\bar{X} = X_{1} + X_{2}$

\item There is $\beta \in [0,1)$ such that $\beta$ is the survival rate of production.
\end{itemize}

\end{frame}

\begin{frame}{Model}

\begin{itemize}
\item Probability of winning the conflict $\Phi_{i} = \Phi_{i}(G_{i},G_{j}) = \frac{G_{i}}{\bar{G}}$

\item $U^{P}_{i} = X_{i}$

\item $U^{C}_{i} = \Phi_{i} \beta \bar{X}$

\item Nash Equilibrium? Conflict is always a NE!

\item If $\beta \geq 0.5$ no stable NE in peace. If $\beta <0.5$, a stable NE in unarmed peace can exist, depending on distribution of resources.
\end{itemize}
\end{frame}

\begin{frame}{Experiment}
\begin{itemize}
\item Measure trust and trustworthinesse, loss aversion and preferences toward preference and risk.

\item Conflict game: Seven rounds, never repeat opponent
\begin{enumerate}
\item Players decide the action(conflict or not) and after this decide how to allocate their endowments.
\item Tell the results.
\end{enumerate}

\item In the case of conflict, $\beta = 0.45$

\item Two different treatment: 

\begin{enumerate} 
\item Equal (100 tokens for each player) 
\item Unequal (120 tokens x 80 tokens)

\end{enumerate}
\end{itemize}
\end{frame}

\begin{frame}
\textbf{Hypothesis 1} Compared to a situation where the distribution of resources is unequal, we expect players to start fewer conflicts when
resources are distributed equally among them.


\textbf{Hypothesis 2}. Compared to a situation where the distribution of resources is unequal, we expect players to invest fewer resources in
arms when resources are distributed equally among them
\end{frame}

\begin{frame}
\includegraphics[scale=0.65]{fig1 - attack.png} 
\end{frame}

\begin{frame}
\begin{figure}[hbtp]
\centering
\includegraphics[scale=0.50]{fig1 - inv.png}

\end{figure}

\end{frame}

\begin{frame}
\begin{figure}[hbtp]

\centering
\includegraphics[scale=0.35]{table2.png}
\end{figure}

\end{frame}

\begin{frame}

\textbf{Result 1:} The probability of an attack is lower when resources are distributed equally among players compared to a situation where the
distribution of resources is unequal. 

\textbf{Result 2:} The average investment in arms is lower when resources are distributed equally among players compared to a situation where the
distribution of resources is unequal

\end{frame}

\begin{frame}
\centering
\includegraphics[scale=0.45]{fig2.png} 
\end{frame}

\begin{frame}
\includegraphics[scale=0.33]{fig3.png} 
\end{frame}

\begin{frame}
\includegraphics[scale=0.65]{fig 4 - attack.png} 
\end{frame}

\begin{frame}
\includegraphics[scale=0.60]{fig 4 - inv.png} 
\end{frame}

\begin{frame}
\includegraphics[scale=0.30]{table3.png} 
\end{frame}


\begin{frame}
\centering
\includegraphics[scale=0.5]{fig 5 - trust.png} 
\end{frame}

\begin{frame}
\begin{figure}[hbtp]
\centering
\includegraphics[scale=0.5]{fig 5 - risk.png}
\end{figure}

\end{frame}

\begin{frame}
\begin{figure}[hbtp]
\centering
\includegraphics[scale=0.55]{fig 5 - loss.png}
\end{figure}

\end{frame}

\begin{frame}{Conclusions}

\textbf{Equal Resources:}
\begin{itemize}
\item Less likely to attack, reduced in arming. Theoretical predicted.

\item Do not find a significant increase of unarmed peace 
\end{itemize}

\textbf{Unequal resources:}
\begin{itemize}

\item Unarmed peace comparable.

\item Disadvantaged players frequently choose peace and show similar tendencies to attack, yet invest less into arming than to
advantaged players. This unexpected finding warrants additional investigation into decisions in conflicts under inequality.

\end{itemize}

\end{frame}

\end{document}